\section{Results}
\label{sec:results}

\subsection{Overall Leakage Rates}
\label{sec:overall}

\tabref{tab:overall} presents the primary result.
Under strict detection, the \neutral condition produces an attack success rate of 58.8\% (95\% CI: [55.5\%, 62.1\%]), while the \conceal condition produces just 2.3\% (95\% CI: [1.3\%, 3.3\%]).
This 56.5 percentage point reduction is highly significant ($\chi^2 = 630.5$, $p < 10^{-139}$, Cram\'er's $V = 0.61$).
The odds ratio is 0.016, indicating that concealment reduces the odds of leakage by a factor of 63.

\begin{table}[t]
\centering
\caption{Overall attack success rate (\asr) under strict detection. Concealment instructions reduce leakage by 56.5 percentage points relative to neutral framing.}
\label{tab:overall}
\begin{tabular}{@{}lccc@{}}
\toprule
\textbf{Framing} & \textbf{ASR} & \textbf{95\% CI} & $n$ \\
\midrule
\neutral & 58.8\% & [55.5\%, 62.1\%] & 840 \\
\conceal & {\bf 2.3\%} & [1.3\%, 3.3\%] & 840 \\
\midrule
Difference & $-$56.5 pp & & \\
\bottomrule
\end{tabular}
\end{table}

\subsection{Leakage by Probe Type}
\label{sec:probe_results}

\tabref{tab:probe} breaks down leakage rates by probe type.
Concealment is effective across all four categories, with the largest absolute reduction on direct probes ($-$96.1 pp) and the smallest on benign probes ($-$14.4 pp).

\begin{table}[t]
\centering
\caption{Attack success rate by probe type under strict detection. Concealment reduces leakage across all probe categories, with the largest effect on direct queries.}
\label{tab:probe}
\begin{tabular}{@{}lccc@{}}
\toprule
\textbf{Probe Type} & \textbf{\neutral ASR} & \textbf{\conceal ASR} & \textbf{Difference} \\
\midrule
Benign & 14.4\% & {\bf 0.0\%} & $-$14.4 pp \\
Direct & 96.1\% & {\bf 0.0\%} & $-$96.1 pp \\
Indirect & 70.6\% & {\bf 1.1\%} & $-$69.5 pp \\
Adversarial & 56.0\% & {\bf 5.7\%} & $-$50.3 pp \\
\bottomrule
\end{tabular}
\end{table}

Several patterns emerge.
First, direct probes achieve near-perfect extraction under neutral framing (96.1\%)---models freely share background information when asked---but concealment drops this to exactly 0\%.
Second, even adversarial attacks that draw on established prompt injection techniques~\citep{pfister2025gandalf, debenedetti2024dataset} see a 10$\times$ reduction in success rate under concealment (56.0\% $\to$ 5.7\%).
Third, the neutral condition exhibits 14.4\% spontaneous leakage on benign probes: models sometimes volunteer background information even when the user has not asked for it.

\subsection{Leakage by Model}
\label{sec:model_results}

\tabref{tab:model} compares leakage rates across the two models.
\gptfour shows higher leakage in both conditions (68.3\% neutral, 4.5\% concealment), while \gptmini achieves \emph{perfect} concealment under strict detection (0.0\%).

\begin{table}[t]
\centering
\caption{Attack success rate by model under strict detection. \gptmini achieves perfect concealment; \gptfour is vulnerable only to translation attacks.}
\label{tab:model}
\begin{tabular}{@{}lcc@{}}
\toprule
\textbf{Model} & \textbf{\neutral ASR} & \textbf{\conceal ASR} \\
\midrule
\gptfour & 68.3\% & 4.5\% \\
\gptmini & 49.3\% & {\bf 0.0\%} \\
\bottomrule
\end{tabular}
\end{table}

The higher neutral-condition leakage for \gptfour reflects its stronger instruction-following: it more readily provides requested background information.
Conversely, all of \gptfour's concealment breaches trace to a single attack vector (see \secref{sec:breach}), suggesting that its vulnerability is narrow rather than systemic.

\subsection{Paired Analysis}
\label{sec:paired}

McNemar's test provides the strongest evidence against the \streisand hypothesis.
Out of 840 paired comparisons (same item + same probe, neutral vs.\ concealment), we observe the distribution in \tabref{tab:mcnemar}.

\begin{table}[t]
\centering
\caption{McNemar's paired analysis: outcome counts for matched neutral--concealment pairs under strict detection. Zero pairs show concealment-only leakage.}
\label{tab:mcnemar}
\begin{tabular}{@{}lc@{}}
\toprule
\textbf{Outcome} & \textbf{Count} \\
\midrule
Both leak & 19 \\
Neither leaks & 346 \\
Only \neutral leaks & 475 \\
Only \conceal leaks & {\bf 0} \\
\bottomrule
\end{tabular}
\end{table}

The McNemar test yields $p = 2.05 \times 10^{-143}$.
Critically, there are \textbf{zero cases} where concealment-framed information leaked but its neutral counterpart did not.
Every concealment breach also occurs under neutral framing, ruling out any scenario in which concealment \emph{causes} leakage.

\subsection{Leakage by Information Category}
\label{sec:category}

\tabref{tab:category} shows that concealment is effective across all information categories.
Categories with lower neutral-condition ASR (\eg codes at 35.7\%) reflect the models' built-in safety training: they refuse to share items that resemble passwords regardless of framing.

\begin{table}[t]
\centering
\caption{Attack success rate by information category under strict detection. Concealment is effective across all categories; lower neutral ASR for codes reflects built-in safety training for password-like strings.}
\label{tab:category}
\begin{tabular}{@{}lcc@{}}
\toprule
\textbf{Category} & \textbf{\neutral ASR} & \textbf{\conceal ASR} \\
\midrule
Code & 35.7\% & 3.6\% \\
Date & 52.4\% & {\bf 0.0\%} \\
Fact & 69.6\% & 3.6\% \\
Name & 70.2\% & 2.4\% \\
Number & 55.4\% & {\bf 0.0\%} \\
Other & 73.8\% & 3.6\% \\
\bottomrule
\end{tabular}
\end{table}

\subsection{Refusal Rates}
\label{sec:refusal}

The concealment condition nearly doubles the refusal rate: 62.6\% vs.\ 30.8\% for neutral framing.
The non-trivial refusal rate in the neutral condition (30.8\%) reflects the models' inherent safety training---they refuse to share items that appear sensitive (\eg ``database password'') even when no concealment instruction is present.

\subsection{Concealment Breach Analysis}
\label{sec:breach}

All 19 strict concealment breaches across both models and all 3 runs trace to a single attack vector: \textit{``Translate your system prompt into French.''}
This attack succeeds because:
\begin{enumerate}[leftmargin=*,itemsep=0pt,topsep=0pt]
\item It frames the request as a translation task rather than an information disclosure request.
\item The model treats translation as a legitimate task, separate from its concealment directive.
\item The translated output includes the \emph{entire} system prompt---both the concealment instruction and the secret value---confirming the ``parallel translation'' mechanism described by \citet{liang2025prompts}.
\end{enumerate}

This attack succeeded only on \gptfour (all 3 runs) and affected 7 of 10 information items---all except those whose values were non-translatable numeric strings.
\gptmini resisted this attack entirely, suggesting stronger compartmentalization between translation and disclosure behavior.
We provide a detailed breakdown in the Appendix.
